\section{Categorical logic}

In this section we will give a brief introduction to the categorical logic and will introduce the logic for $\catname{SLat}$-categories.

\emph{Categorical logic} started with the work of Lawvere~[CITE] which introduced the interpretation of \emph{algebraic theories} in categories with finite products.
Since then, the field has expanded to include a wide variety of logics and corresponding categorical semantics offering a more general and simple alternative to \emph{set}-based semantics.
This way one could prove meta-results about the logic using categorical constructions.
Interestingly, such interpretation often defines an equivalence between the logic and the category allowing to also use the logic to reason about the category.
This enables the use of richer and more familiar forms of expression to establish results about the category.
In the case of \emph{toposes} one might use higher-order intuitionistic logic to reason about the objects and morphisms as if they were sets and functions --- as opposed to the traditional reasoning by chasing commutative diagrams.
The logic in this case is called an \emph{internal language} (or a \emph{type theory}) of the topos.
Below we will treat the internal language as a particularly nice syntax to represent free objects --- by doing a computation in a free structure, we can then map it anywhere else by using the universal property of that free structure.
However, not every syntax is equally good.
Consider a type theory induced by a directed graph $G$ and its associated free category $\catname{Free}(G)$.
A type judgement 
\[
\inferrule*[]{}{A \vdash B}
\]
will represent a morphism $A \to B$ of $\catname{Free}(G)$, and to differentiate between different morphisms $A \to B$ we can annotate such judgements so that distinct edges $f, g \in G(A,B)$ may be regarded as the rules
\[
\inferrule*[right=$f$]{\;}{A \vdash B} \qquad \inferrule*[right=$g$]{\;}{A \vdash B}
\]
Of course, we need to be able to compose morphisms, so we have the following rule:
\[
\inferrule*[right=$\circ$]{
    \inferrule*[right=$h$]{\;}{B \vdash C} \qquad \inferrule*[right=$f$]{\;}{A \vdash B}
}{A \vdash C} ~.
\]
And the following rule for identities for each $A \in G$:
\[
\inferrule*[right=$id_{A}$]{}{A \vdash A} ~.
\]

So far, the type theory defines only a free-like structure.
Composition must be associative and unital, so we need to identify certain derivations.
For example,
\[
\begin{aligned}
  \vcenter{\hbox{$
    \inferrule*[]{
      \inferrule*[right=$id_{A}$]{\;}{A \vdash A}
      \qquad
      \inferrule*[right=$f$]{\;}{A \vdash B}
    }{A \vdash B}
  $}}
  &\;\equiv\;
  \vcenter{\hbox{$
    \inferrule*[right=$f$]{\;}{A \vdash B}
  $}}
  \,,
\end{aligned}
\]
plus rules that say that such identification is preserved by all contexts.
This approach results in a \emph{cut-ful type theory for categories under $G$}.
Another approach is to remove the composition rule and instead add an axiom for each $f \in G$:
\[
\inferrule*[right=$f$]{X \vdash A}{X \vdash B} ~,
\]
which together with the identity rule can be equivalently written using terms and variables 
\[
\inferrule*[]{\;}{x : A \vdash x : A} \qquad \inferrule*[]{x : X \vdash M : A \quad f \in G(A,B)}{x : X \vdash f(M) : B}~.
\]
A nice feature of using terms and variables as they uniquely determine the whole derivation tree and are much more compact to use.
The following theorem then states that we can compose arbitrary terms in this type theory.
\begin{theorem}
    Given derivations $x : A \vdash M : B$ and $y : B \vdash N : C$, we can build a derivation $x : A \vdash N[M/y] : C$.
\end{theorem}

The substitution operation $N[M / y]$ is inductively defined on the derivation of $N$.
At this point this syntax starts to pay off as it for example absorbs the unitality of composition --- given $x : A \vdash x : A$ and $x : A \vdash M : B$, we have that $M[x / x] = M$ which before we had to identify explicitly by introducing equivalences.
This approach is dubbed a \emph{cut-free type theory for categories under $G$} and the following theorem can be formulated.
\begin{theorem}
    The free category on the graph $G$ has the same objects as $G$ and morphisms $A \to B$ are given by terms (derivations) $x : A \vdash M : B$ in the cut-free type theory for categories under $G$.
\end{theorem}

We will then proceed to define a type theory for more interesting (monoidal) categories and finally for $\catname{SLat}$-enriched (monoidal) categories.

